\documentclass[11pt,reqno]{amsart}

% ================================================================
% Packages
% ================================================================
\usepackage{amsmath,amssymb,amsthm}
\usepackage{mathrsfs}
\usepackage{enumerate}
\usepackage[shortlabels]{enumitem}
\usepackage{booktabs}
\usepackage{array}
\usepackage{microtype}
\usepackage{tikz-cd}
\usepackage[dvipsnames]{xcolor}
\usepackage[colorlinks=true,
  linkcolor=blue!60!black,
  citecolor=green!40!black,
  urlcolor=blue!60!black]{hyperref}

% ================================================================
% Theorem environments
% ================================================================
\newtheorem{theorem}{Theorem}[section]
\newtheorem{proposition}[theorem]{Proposition}
\newtheorem{lemma}[theorem]{Lemma}
\newtheorem{corollary}[theorem]{Corollary}
\newtheorem{conjecture}[theorem]{Conjecture}
\theoremstyle{definition}
\newtheorem{definition}[theorem]{Definition}
\newtheorem{construction}[theorem]{Construction}
\newtheorem{example}[theorem]{Example}
\newtheorem{warning}[theorem]{Warning}
\theoremstyle{remark}
\newtheorem{remark}[theorem]{Remark}
\newtheorem{notation}[theorem]{Notation}

% ================================================================
% Macros
% ================================================================
\newcommand{\cA}{\mathcal{A}}
\newcommand{\cB}{\mathcal{B}}
\newcommand{\cC}{\mathcal{C}}
\newcommand{\cD}{\mathcal{D}}
\newcommand{\cF}{\mathcal{F}}
\newcommand{\cG}{\mathcal{G}}
\newcommand{\cH}{\mathcal{H}}
\newcommand{\cL}{\mathcal{L}}
\newcommand{\cM}{\mathcal{M}}
\newcommand{\cO}{\mathcal{O}}
\newcommand{\cP}{\mathcal{P}}
\newcommand{\cQ}{\mathcal{Q}}
\newcommand{\cW}{\mathcal{W}}
\newcommand{\cZ}{\mathcal{Z}}

\newcommand{\CC}{\mathbb{C}}
\newcommand{\RR}{\mathbb{R}}
\newcommand{\ZZ}{\mathbb{Z}}
\newcommand{\PP}{\mathbb{P}}
\providecommand{\Bbbk}{\mathbf{k}}

\newcommand{\fg}{\mathfrak{g}}
\newcommand{\fh}{\mathfrak{h}}
\newcommand{\fn}{\mathfrak{n}}
\newcommand{\fb}{\mathfrak{b}}

\newcommand{\Ran}{\mathrm{Ran}}
\newcommand{\MC}{\mathrm{MC}}
\newcommand{\HH}{\mathrm{HH}}
\newcommand{\CE}{\mathrm{CE}}
\newcommand{\PVA}{\mathrm{PVA}}
\newcommand{\Sym}{\mathrm{Sym}}
\newcommand{\id}{\mathrm{id}}
\newcommand{\op}{\mathrm{op}}
\newcommand{\barB}{\overline{B}}
\newcommand{\barA}{\bar{\cA}}

\newcommand{\Ainf}{A_{\infty}}
\newcommand{\Linf}{L_{\infty}}
\newcommand{\Pinf}{P_{\infty}}
\newcommand{\Eone}{\mathsf{E}_{1}}
\newcommand{\Etwo}{\mathsf{E}_{2}}
\newcommand{\Ethree}{\mathsf{E}_{3}}
\newcommand{\En}{\mathsf{E}_{n}}
\newcommand{\Einf}{\mathsf{E}_{\infty}}

\newcommand{\SC}{\mathsf{SC}}
\newcommand{\SCchtop}{\mathsf{SC}^{\mathrm{ch,top}}}

\newcommand{\Ass}{\mathsf{Ass}}
\newcommand{\Com}{\mathsf{Com}}
\newcommand{\Lie}{\mathsf{Lie}}

\DeclareMathOperator{\Hom}{Hom}
\DeclareMathOperator{\RHom}{RHom}
\DeclareMathOperator{\Res}{Res}
\DeclareMathOperator{\Aut}{Aut}
\DeclareMathOperator{\End}{End}
\DeclareMathOperator{\Conf}{Conf}
\DeclareMathOperator{\FM}{FM}
\DeclareMathOperator{\Spec}{Spec}
\DeclareMathOperator{\av}{av}
\DeclareMathOperator{\rk}{rk}
\DeclareMathOperator{\Tr}{Tr}
\DeclareMathOperator{\ad}{ad}
\DeclareMathOperator{\colim}{colim}

\numberwithin{equation}{section}

% Claim status macros (standalone: no-op definitions)
\providecommand{\ClaimStatusConjectured}{}
\providecommand{\ClaimStatusProvedHere}{}
\providecommand{\ClaimStatusProvedElsewhere}{}
\providecommand{\ClaimStatusHeuristic}{}

% ================================================================
\begin{document}

\title[SC$^{\mathrm{ch,top}}$ and PVA descent]
{The Swiss-cheese operad $\SCchtop$ and PVA descent}

\author{Raeez Lorgat}
\address{Perimeter Institute for Theoretical Physics,
  Waterloo, ON N2L 2Y5, Canada \\
  Department of Physics, University of Waterloo,
  Waterloo, ON N2L 3G1, Canada}
\email{rlorgat@perimeterinstitute.ca}

\date{\today}

\begin{abstract}
A chiral algebra $\cA$ on a smooth curve $X$ determines a bar
complex $\barB(\cA) = T^c(s^{-1}\barA)$ that is an $\Eone$
chiral coassociative coalgebra: the differential encodes
holomorphic factorisation on $\FM_k(\CC)$, the deconcatenation
coproduct encodes topological factorisation on
$\Conf_k^{<}(\RR)$.  The two-coloured operad governing the
interaction of these two structures on the \emph{derived chiral
centre} $C^\bullet_{\mathrm{ch}}(\cA,\cA)$ and the boundary
algebra~$\cA$ is the Swiss-cheese operad
$\SCchtop$, with closed colour $\FM_k(\CC)$, open colour
$\Eone(m) = \Conf_m^{<}(\RR)$, and empty open-to-closed
component (directionality).

We establish three results.
First, $\SCchtop$ is homotopy-Koszul, via Kontsevich formality
for the closed colour composed with Livernet's Koszulity of the
classical Swiss-cheese operad $\SC$ and bar-cobar transfer.
Second, the Koszul dual cooperad satisfies
$(\SCchtop)^! = (\Lie^c,\,\Ass^c,\,\text{shuffle-mixed})$;
in particular $\SCchtop$ is \emph{not} Koszul self-dual
($\dim\SCchtop_{\mathrm{cl}}(n) = 1$ versus
$\dim(\SCchtop)^!_{\mathrm{cl}}(n) = (n-1)!$).
Third, the cogenerator projections of the bar differential
produce $\Ainf$ operations $\{m_k\}_{k \geq 1}$ whose Stasheff
relations are Stokes' theorem on $\FM_n(\CC)$; on cohomology
these descend to the five axioms of a Poisson vertex algebra.
The descent is computed in three cases:
the Heisenberg algebra (formal, class~G, shadow depth zero),
affine Kac--Moody at level~$k$ (Lie-transverse, class~L, shadow
depth one), and $\cW_3$ (infinite tower, class~M).
\end{abstract}

\maketitle
\tableofcontents


%%%%%%%%%%%%%%%%%%%%%%%%%%%%%%%%%%%%%%%%%%%%%%%%%%%%%%%%%%%%%
\section{Introduction: two colours, one operad}
\label{sec:introduction}
%%%%%%%%%%%%%%%%%%%%%%%%%%%%%%%%%%%%%%%%%%%%%%%%%%%%%%%%%%%%%

\subsection{The problem}

The bar complex of a chiral algebra carries two independent
structures.  The \emph{differential} $D_\cA$ is a coderivation
of the cofree tensor coalgebra $T^c(s^{-1}\barA)$; it encodes
the OPE residues extracted from collisions on the
Fulton--MacPherson spaces $\FM_k(\CC)$.  The
\emph{deconcatenation coproduct} $\Delta$ cuts the tensor word
at each gap:
\begin{equation}\label{eq:intro-deconc}
  \Delta[s^{-1}a_1 \mid \cdots \mid s^{-1}a_n]
  \;=\;
  \sum_{i=0}^{n}
  [s^{-1}a_1 \mid \cdots \mid s^{-1}a_i]
  \;\otimes\;
  [s^{-1}a_{i+1} \mid \cdots \mid s^{-1}a_n].
\end{equation}
The compatibility between these two structures is the
coderivation property:
\begin{equation}\label{eq:intro-coderivation}
  \Delta \circ D_\cA
  \;=\;
  (D_\cA \otimes \id + \id \otimes D_\cA)
  \circ \Delta.
\end{equation}
A differential and a coassociative coproduct satisfying
\eqref{eq:intro-coderivation} make $(\barB(\cA), D_\cA, \Delta)$
into a dg coassociative coalgebra: an $\Eone$ coalgebra, with
holomorphic collisions on $\FM_k(\CC)$ as the differential and
topological separations on $\Conf_k^{<}(\RR)$ as the coproduct.

The question is: \emph{what operad governs the structures
that emerge from this $\Eone$ engine?}

\subsection{The answer}

The bar complex itself does not carry the answer.  It is the
input, not the output.  The output is the \emph{chiral derived
centre}: the chiral Hochschild cochain complex
$C^\bullet_{\mathrm{ch}}(\cA,\cA)$, defined via the chiral
endomorphism operad $\End^{\mathrm{ch}}_\cA$ with spectral
parameters from $\FM_k(\CC)$.  The pair
\begin{equation}\label{eq:intro-SC-pair}
  \bigl(C^\bullet_{\mathrm{ch}}(\cA,\cA),\; \cA\bigr)
  \;=\;
  (\text{bulk},\; \text{boundary})
\end{equation}
is the datum governed by the two-coloured Swiss-cheese operad
$\SCchtop$.  The closed colour is the bulk; the open colour is
the boundary.  Bulk acts on boundary; boundary does not
act on bulk.  This directionality is the empty open-to-closed
component of $\SCchtop$.

Three properties hold.

\begin{enumerate}[label=(\roman*)]
\item \textbf{Homotopy-Koszulity.}  The operad $\SCchtop$ is
  homotopy-Koszul (Theorem~\ref{thm:homotopy-koszul-SC}): the
  bar-cobar adjunction
  $\Omega\mathbf{B}(\SCchtop) \xrightarrow{\sim} \SCchtop$
  is a quasi-isomorphism.  This follows from Livernet's
  Koszulity of the classical Swiss-cheese operad~$\SC$
  \cite{Liv06}, composed with Kontsevich formality for
  the closed colour, and transferred to $\SCchtop$ via
  bar-cobar comparison.

\item \textbf{The Koszul dual cooperad.}  The Koszul dual
  satisfies
  $(\SCchtop)^! = (\Lie^c,\, \Ass^c,\, \text{shuffle-mixed})$
  (Theorem~\ref{thm:SC-koszul-dual}).  The closed colour
  exchanges $\Com \leftrightarrow \Lie$; the open colour
  preserves $\Ass$ (which is self-dual); the mixed sector
  acquires a shuffle structure.  In particular, $\SCchtop$ is
  \emph{not} Koszul self-dual: the closed-colour component
  dimensions differ ($\dim\SCchtop_{\mathrm{cl}}(n) = 1$ versus
  $\dim(\SCchtop)^!_{\mathrm{cl}}(n) = (n-1)!$).

\item \textbf{PVA descent.}  The cogenerator projections of
  $D_\cA$ yield $\Ainf$ operations $\{m_k\}_{k \geq 1}$
  satisfying the Stasheff relations, each relation being
  Stokes' theorem on $\FM_n(\CC)$.  On cohomology, the
  operations $m_k$ descend to a commutative product and a
  $\lambda$-bracket satisfying the five axioms of a Poisson
  vertex algebra: sesquilinearity, skew-symmetry, Jacobi,
  Leibniz, and commutativity
  (Theorem~\ref{thm:pva-descent}).
\end{enumerate}

\begin{warning}[The bar complex is not an
  $\SCchtop$-coalgebra]\label{warn:bar-not-SC}
The bar complex $\barB(\cA) = T^c(s^{-1}\barA)$ is an $\Eone$
chiral coassociative coalgebra.  It carries a differential
(from the chiral product) and a coproduct
(from deconcatenation).  These are two structures on a
\emph{single} coalgebra.  The $\SCchtop$ structure is
\emph{not} carried by $\barB(\cA)$; it is carried by the
pair $(C^\bullet_{\mathrm{ch}}(\cA,\cA),\, \cA)$, where
$C^\bullet_{\mathrm{ch}}(\cA,\cA)$ is computed \emph{using}
the bar complex as a resolution.  The bar complex is the
$\Eone$ engine; the derived centre pair is the $\SCchtop$
output.
\end{warning}

\subsection{Context and conventions}

We work over a field $\Bbbk$ of characteristic zero.  Grading
is cohomological: $|d| = +1$.  Desuspension lowers degree:
$|s^{-1}v| = |v| - 1$.  The bar complex uses the augmentation
ideal: $\barB(\cA) = T^c(s^{-1}\barA)$ where
$\barA = \ker(\varepsilon)$.  The Fulton--MacPherson space
$\FM_n(X)$ is the iterated real blowup of $X^n$ along all
diagonal strata.  Configuration spaces of ordered points on
$\RR$ are denoted $\Conf_n^{<}(\RR) = \{t_1 < \cdots < t_n\}$.

Throughout, $\cA$ denotes a chiral algebra on a smooth algebraic
curve $X$ over $\Bbbk$ in the sense of Beilinson--Drinfeld
\cite{BD04}.  A vertex algebra is a chiral algebra on the
formal disk $D = \Spec\Bbbk[[z]]$.

\subsection{Relation to other work}

The classical Swiss-cheese operad $\SC$ (with $\Etwo$ closed
colour and $\Eone$ open colour) was introduced by Voronov
\cite{Vor99}.  Koszulity of $\SC$ was proved by Livernet
\cite{Liv06}.  The deformation-theoretic role of $\SC$ in
open-closed homotopy algebras was developed by Kajiura--Stasheff
\cite{KS06}.  The passage from the classical $\SC$ to the chiral
variant $\SCchtop$ replaces the $\Etwo$ closed colour by
$\FM_k(\CC)$-spaces carrying holomorphic structure; this
replacement is the content of the Beilinson--Drinfeld
factorisation framework \cite{BD04}.  The relationship between
$\Ainf$ structures and PVA axioms has been studied from the
$\lambda$-bracket perspective by De Sole and Kac \cite{DSK06},
and from the factorisation perspective by Francis and Gaitsgory
\cite{FG12}.  The present paper makes the operadic architecture
of this descent explicit.


%%%%%%%%%%%%%%%%%%%%%%%%%%%%%%%%%%%%%%%%%%%%%%%%%%%%%%%%%%%%%
\section{The $\SCchtop$ operad}
\label{sec:SC-operad}
%%%%%%%%%%%%%%%%%%%%%%%%%%%%%%%%%%%%%%%%%%%%%%%%%%%%%%%%%%%%%

\subsection{Configuration spaces}

\begin{definition}[Fulton--MacPherson space]
\label{def:FM-space}
For a smooth manifold $M$ and $n \geq 1$, the
\emph{Fulton--MacPherson space} $\FM_n(M)$ is the real-oriented
blowup of $M^n$ along all diagonal strata
$\Delta_S = \{(x_i)_{i \in I} : x_i = x_j \text{ for all }
i,j \in S\}$, for all subsets $S \subseteq \{1,\ldots,n\}$ with
$|S| \geq 2$, performed in order of increasing $|S|$.  Points
of $\FM_n(M)$ are configurations of $n$ labelled points in $M$,
together with ``infinitesimal data'' recording the relative rates
and directions of approach to collision.
\end{definition}

\begin{definition}[Ordered configuration space]
\label{def:ordered-conf}
For $n \geq 1$, define
\[
  \Conf_n^{<}(\RR)
  \;=\;
  \{(t_1, \ldots, t_n) \in \RR^n : t_1 < t_2 < \cdots < t_n\}.
\]
This is contractible: $\Conf_n^{<}(\RR) \cong \RR^{n-1}_{>0}$
via the change of variables $u_i = t_{i+1} - t_i$.
\end{definition}

\subsection{The two-coloured operad}

\begin{definition}[The Swiss-cheese operad $\SCchtop$]
\label{def:SC-chtop}
The operad $\SCchtop$ is a two-coloured operad with colours
$\{\mathsf{cl}, \mathsf{op}\}$ (closed and open), defined by
the following operation spaces:
\begin{enumerate}[label=(\roman*)]
\item \textbf{Closed-to-closed.}  For $k \geq 1$ closed inputs
  and one closed output:
  \begin{equation}\label{eq:SC-cl-cl}
    \SCchtop(\mathsf{cl}^k;\, \mathsf{cl})
    \;=\;
    \FM_k(\CC).
  \end{equation}
  Composition is given by the operadic structure maps of the
  Fulton--MacPherson operad: for a decomposition
  $\{1,\ldots,k\} = S_1 \sqcup \cdots \sqcup S_r$, the
  composition map
  \[
    \gamma \colon
    \FM_r(\CC) \times \FM_{|S_1|}(\CC) \times \cdots \times
    \FM_{|S_r|}(\CC)
    \;\longrightarrow\;
    \FM_k(\CC)
  \]
  inserts the inner configurations into the outer one at the
  collision points.

\item \textbf{Open-to-open.}  For $m \geq 1$ open inputs and
  one open output:
  \begin{equation}\label{eq:SC-op-op}
    \SCchtop(\mathsf{op}^m;\, \mathsf{op})
    \;=\;
    \Eone(m)
    \;=\;
    \Conf_m^{<}(\RR).
  \end{equation}
  Composition is by interval insertion: the $\Eone$ operadic
  structure on ordered configurations.  Since
  $\Conf_m^{<}(\RR)$ is contractible, the open-to-open
  component is equivalent to the associative operad $\Ass$
  up to homotopy.

\item \textbf{Mixed (closed and open inputs, open output).}
  For $k \geq 0$ closed inputs and $m \geq 1$ open inputs
  with one open output:
  \begin{equation}\label{eq:SC-mix}
    \SCchtop(\mathsf{cl}^k, \mathsf{op}^m;\, \mathsf{op})
    \;=\;
    \FM_k(\CC) \times \Eone(m).
  \end{equation}
  The composition is the product of FM-substitution on the
  closed factor and interval-insertion on the open factor.

\item \textbf{Directionality (open-to-closed is empty).}
  For any input profile containing at least one open colour
  and producing a closed output:
  \begin{equation}\label{eq:SC-direction}
    \SCchtop(\ldots, \mathsf{op}, \ldots;\, \mathsf{cl})
    \;=\;
    \varnothing.
  \end{equation}
  Bulk restricts to boundary; boundary does not act on bulk.
\end{enumerate}
\end{definition}

\begin{remark}[Directionality and the Springer pattern]
\label{rem:directionality}
The empty open-to-closed component is a structural feature,
not a convention.  In the Beilinson--Drinfeld framework, the
factorisation $D$-module on $\Ran(X)$ (the closed sector)
restricts to the boundary via $\iota^!$, but the reverse
extension $\iota_!$ produces zero $!$-tensor product with the
bulk $D$-module by support considerations: the boundary
$X \times \{0\}$ has codimension one in $X \times \RR$, and
the $!$-extension carries no sections transverse to the
boundary.
\end{remark}

\begin{remark}[Comparison with the classical $\SC$]
\label{rem:classical-SC}
Voronov's classical Swiss-cheese operad $\SC$ has $\Etwo$ closed
colour (configurations of little disks in a disk) and $\Eone$
open colour (configurations of little intervals in an interval).
The operad $\SCchtop$ replaces the $\Etwo$ closed colour by
$\FM_k(\CC)$-spaces that carry holomorphic structure: the complex
structure on $\CC$ is visible in the boundary strata of
$\FM_k(\CC)$ and enters the bar differential through the OPE
residues.  On cohomology, the Kontsevich formality theorem
provides a quasi-isomorphism from $\FM_k(\CC)$-chains to the
$\Etwo$-operad; composing with the projection from $\SC$ to the
homology operad $(\Com, \Ass)$ recovers the classical descent to
PVA\@.  At the chain level, the holomorphic structure is load-bearing:
it determines the spectral parameters of the $\Ainf$ operations.
\end{remark}

\begin{definition}[$\SCchtop$-algebra]
\label{def:SC-algebra}
An \emph{$\SCchtop$-algebra} is a pair $(B, A)$ where:
\begin{enumerate}[label=(\roman*)]
\item $A$ is an $\Eone$-algebra (the open/boundary algebra);
\item $B$ is an algebra over the FM-operad
  $\{\FM_k(\CC)\}_{k \geq 1}$ (the closed/bulk algebra);
\item the mixed operations
  $\FM_k(\CC) \times \Eone(m) \to
  \Hom(B^{\otimes k} \otimes A^{\otimes m},\, A)$
  define a module structure of $B$ over $A$, compatible with
  all composition maps;
\item the open-to-closed maps are zero.
\end{enumerate}
\end{definition}

\begin{example}[The chiral derived centre pair]
\label{ex:derived-centre-pair}
For a chiral algebra $\cA$ on $X$, the chiral Hochschild cochain
complex
\[
  C^\bullet_{\mathrm{ch}}(\cA,\cA)
  \;:=\;
  \Hom_{\barB(\cA)}\bigl(\barB(\cA),\, \cA\bigr)
\]
(the convolution $\Hom$ from the bar resolution into $\cA$) carries
the structure of an algebra over $\{\FM_k(\CC)\}_{k \geq 1}$
via the chiral brace operations.  The pair
$(C^\bullet_{\mathrm{ch}}(\cA,\cA),\, \cA)$ is an
$\SCchtop$-algebra.  The closed colour is the bulk; the open
colour is the boundary.  The bar complex $\barB(\cA)$ is used
\emph{in the construction} of $C^\bullet_{\mathrm{ch}}(\cA,\cA)$
but does not itself carry the $\SCchtop$ structure.
\end{example}


%%%%%%%%%%%%%%%%%%%%%%%%%%%%%%%%%%%%%%%%%%%%%%%%%%%%%%%%%%%%%
\section{Homotopy-Koszulity}
\label{sec:homotopy-koszulity}
%%%%%%%%%%%%%%%%%%%%%%%%%%%%%%%%%%%%%%%%%%%%%%%%%%%%%%%%%%%%%

\subsection{Koszulity of the classical Swiss-cheese operad}

Recall that a quadratic operad $\cP$ is \emph{Koszul} if the
natural augmentation map
$\Omega(\cP^!) \xrightarrow{\sim} \cP$ is a quasi-isomorphism,
where $\cP^!$ is the Koszul dual cooperad and $\Omega$ denotes
the cobar functor.  For coloured operads, the Koszul property is
defined colour by colour with compatibility conditions on the
mixed sectors.

\begin{theorem}[Livernet {\cite{Liv06}}]
\label{thm:livernet-koszulity}
\ClaimStatusProvedElsewhere{}
The classical Swiss-cheese operad $\SC$, with $\Etwo$ closed
colour and $\Eone$ open colour, is Koszul.
\end{theorem}

The proof proceeds by establishing a distributive law between
the $\Com$ and $\Ass$ components that governs the mixed sector,
then verifying the rewriting criterion of Dotsenko--Khoroshkin
\cite{DK10} for the resulting coloured quadratic presentation.

\subsection{Formality of the closed colour}

The closed colour of $\SCchtop$ is the FM-operad
$\{\FM_k(\CC)\}_{k \geq 1}$.  Kontsevich's formality theorem
\cite{Kon99,Kon03} provides a quasi-isomorphism of operads:
\begin{equation}\label{eq:kontsevich-formality}
  \{H_\bullet(\FM_k(\CC))\}_{k \geq 1}
  \;\xleftarrow{\;\sim\;}\;
  \{\FM_k(\CC)\}_{k \geq 1},
\end{equation}
where the left-hand side is the homology operad, which is the
$\Etwo$-operad $\{H_\bullet(\Conf_k(\CC))\}$ concentrated in
degree zero.  The formality is over $\RR$; over $\CC$ it holds
unconditionally.  The key input is the propagator on $\FM_2(\CC)$
and the vanishing of all obstructions via the Arnold relation.

\begin{remark}[Formality is for the closed colour, not the
  full operad]\label{rem:formality-scope}
Kontsevich formality applies to the closed colour of $\SCchtop$
(the FM-operad on $\CC$).  It does \emph{not} assert formality
of the full two-coloured operad $\SCchtop$; the mixed sector
carries non-trivial chain-level data that does not simplify under
the formality quasi-isomorphism.  The chain-level mixed data is
precisely the spectral-parameter dependence of the $\Ainf$
operations.
\end{remark}

\subsection{Transfer to $\SCchtop$}

\begin{theorem}[Homotopy-Koszulity of $\SCchtop$]
\label{thm:homotopy-koszul-SC}
\ClaimStatusProvedHere{}
The operad $\SCchtop$ is homotopy-Koszul: the canonical map
\begin{equation}\label{eq:bar-cobar-SC}
  \Omega\mathbf{B}(\SCchtop)
  \;\xrightarrow{\;\sim\;}\;
  \SCchtop
\end{equation}
is a quasi-isomorphism.  The bar-cobar adjunction on
$\SCchtop$-algebras is a Quillen equivalence.
\end{theorem}

\begin{proof}
The argument has three steps.

\medskip\noindent\textbf{Step~1.}
Livernet's Theorem~\ref{thm:livernet-koszulity} gives
Koszulity of the classical $\SC$: the bar-cobar
resolution $\Omega(\SC^!) \xrightarrow{\sim} \SC$ is a
quasi-isomorphism.

\medskip\noindent\textbf{Step~2.}
Kontsevich formality \eqref{eq:kontsevich-formality} provides
a chain of quasi-isomorphisms connecting the closed colour
of $\SCchtop$ (the FM-operad) to the closed colour of $\SC$
(the $\Etwo$-operad, the homology of FM).  The open colour is
already $\Eone$ in both operads, so no formality argument is
needed there.

\medskip\noindent\textbf{Step~3.}
The bar-cobar transfer principle
(Loday--Vallette \cite{LV12}, Theorem~11.3.3) states: if
$f \colon \cP \xrightarrow{\sim} \cQ$ is a quasi-isomorphism
of dg operads and $\cQ$ is Koszul, then $\cP$ is
homotopy-Koszul; that is, the natural map
$\Omega\mathbf{B}(\cP) \xrightarrow{\sim} \cP$ is a
quasi-isomorphism.  Applying this with $\cP = \SCchtop$ and
$\cQ = \SC$ (connected by the formality quasi-isomorphism on
the closed colour and the identity on the open colour) yields
the result.

The Quillen equivalence follows from the general model category
structure on algebras over a homotopy-Koszul operad
(Hinich \cite{Hin97}, extended to the coloured setting by
Caviglia \cite{Cav14}).
\end{proof}

\begin{remark}[Why homotopy-Koszulity matters]
\label{rem:homotopy-koszul-consequences}
Homotopy-Koszulity of $\SCchtop$ has three consequences:
\begin{enumerate}[label=(\alph*)]
\item The bar construction $\mathbf{B}(\SCchtop)$ is a
  well-defined cooperad with $\partial^2 = 0$; this is the
  prerequisite for the convolution algebra
  $\Hom_\Sigma(\mathbf{B}(\SCchtop),\, \End_{(B,A)})$ to
  carry an $\Linf$-structure and for the MC element $\alpha_T$
  to exist.
\item The cobar resolution $\Omega\mathbf{B}(\SCchtop)$ provides
  a cofibrant replacement of $\SCchtop$; algebras over this
  resolution are $\SCchtop$-algebras up to homotopy.
\item The homotopy transfer theorem applies: given a deformation
  retract of the pair $(B,A)$ onto cohomology $(H^\bullet(B),
  H^\bullet(A))$, the $\SCchtop$-algebra structure transfers to
  a homotopy $\SCchtop$-algebra structure on cohomology.  This
  transfer is the mechanism producing the PVA structure.
\end{enumerate}
\end{remark}


%%%%%%%%%%%%%%%%%%%%%%%%%%%%%%%%%%%%%%%%%%%%%%%%%%%%%%%%%%%%%
\section{The Koszul dual: $(\SCchtop)^!$ is not self-dual}
\label{sec:koszul-dual}
%%%%%%%%%%%%%%%%%%%%%%%%%%%%%%%%%%%%%%%%%%%%%%%%%%%%%%%%%%%%%

\subsection{The Koszul dual cooperad}

\begin{theorem}[Koszul dual of $\SCchtop$]
\label{thm:SC-koszul-dual}
\ClaimStatusProvedHere{}
The Koszul dual cooperad of $\SCchtop$ is
\begin{equation}\label{eq:SC-dual}
  (\SCchtop)^!
  \;=\;
  (\Lie^c,\; \Ass^c,\; \text{shuffle-mixed}),
\end{equation}
with the following structure:
\begin{enumerate}[label=(\roman*)]
\item \textbf{Closed colour.}
  $(\SCchtop)^!_{\mathrm{cl}}(n) = \Lie^c(n)$, the Lie
  cooperad.  The dimension is
  $\dim (\SCchtop)^!_{\mathrm{cl}}(n) = (n-1)!$.
\item \textbf{Open colour.}
  $(\SCchtop)^!_{\mathrm{op}}(m) = \Ass^c(m) = \Bbbk$,
  the associative cooperad, one-dimensional in each degree.
\item \textbf{Mixed sector.}
  $(\SCchtop)^!(\mathsf{cl}^k, \mathsf{op}^m;\, \mathsf{op})$
  carries the shuffle cooperad structure:
  \begin{equation}\label{eq:SC-dual-mixed}
    \dim(\SCchtop)^!(\mathsf{cl}^k, \mathsf{op}^m;\, \mathsf{op})
    \;=\;
    (k-1)! \cdot \binom{k+m}{m}.
  \end{equation}
\end{enumerate}
\end{theorem}

\begin{proof}
The Koszul duality for two-coloured operads
(Loday--Vallette \cite{LV12}, Chapter~8) exchanges each
colour's operad with its Koszul dual and intertwines the mixed
sector through the distributive law.

For the closed colour: $\Com^! = \Lie$ (the classical
Ginzburg--Kapranov duality \cite{GK94}).  The operation spaces
exchange: $\dim\Com(n) = 1$ becomes
$\dim\Lie(n) = (n-1)!$.

For the open colour: $\Ass^! = \Ass$ (associative is
self-dual under Koszul duality).  The operation spaces
are one-dimensional in each degree, and this is preserved.

For the mixed sector: the distributive law of
$\Com$ over $\Ass$ dualises to the shuffle structure.  The
dimension formula \eqref{eq:SC-dual-mixed} follows from
the shuffle product: $(k-1)!$ from the Lie component and
$\binom{k+m}{m}$ from the shuffle of closed and open inputs.
\end{proof}

\begin{warning}[$\SCchtop$ is not Koszul self-dual]
\label{warn:not-self-dual}
The closed-colour dimensions differ:
\begin{center}
\renewcommand{\arraystretch}{1.2}
\begin{tabular}{lcccccc}
  \toprule
  $n$ & 1 & 2 & 3 & 4 & 5 & 6 \\
  \midrule
  $\dim \SCchtop_{\mathrm{cl}}(n) = \dim\Com(n)$
  & 1 & 1 & 1 & 1 & 1 & 1 \\
  $\dim (\SCchtop)^!_{\mathrm{cl}}(n) = \dim\Lie(n)$
  & 1 & 1 & 2 & 6 & 24 & 120 \\
  \bottomrule
\end{tabular}
\end{center}
At $n \geq 3$ these are manifestly different.  The operad
$\SCchtop$ is therefore \emph{not} Koszul self-dual.

What \emph{is} true: the Koszul duality \emph{functor} on
$\SCchtop$-algebras is an involution.  For any Koszul operad
$\cP$, the double dual satisfies $(\cP^!)^! \cong \cP$
(tautologically).  This means the duality functor
$\cA \mapsto \cA^!$ on algebras is an involution; it does
\emph{not} mean the operad is isomorphic to its own dual.
The confusion between ``the functor is an involution'' and
``the operad is self-dual'' is a common source of error.
\end{warning}

\subsection{Algebras over $(\SCchtop)^!$}

\begin{proposition}[Structure of an $(\SCchtop)^!$-algebra]
\label{prop:SC-dual-algebra}
\ClaimStatusProvedHere{}
An algebra over $(\SCchtop)^!$ is a pair $(L, A)$ where:
\begin{enumerate}[label=(\roman*)]
\item $L$ is a Lie algebra (the closed colour, from the
  $\Lie$ component);
\item $A$ is an associative algebra (the open colour, from
  the $\Ass$ component);
\item $L$ acts on $A$ by derivations, through the mixed
  sector;
\item the open-to-closed direction is zero.
\end{enumerate}
The chiral Koszul dual $\cA^!_{\mathrm{line}}$ of a chiral
algebra carries exactly this structure: a Lie bracket
(the Sklyanin bracket from the closed colour) and an
associative product (the Yangian product from the open colour),
with the Lie bracket acting as derivations.
\end{proposition}

\begin{proof}
This is a direct consequence of
Theorem~\ref{thm:SC-koszul-dual}.  The $\Lie$ cooperad
dualises to the $\Lie$ operad on algebras; the $\Ass$ cooperad
dualises to the $\Ass$ operad on algebras; the shuffle mixed
sector dualises to the derivation condition.
\end{proof}

\begin{remark}[The asymmetry principle]
\label{rem:asymmetry}
The Koszul dual exchanges $\Com$ (one-dimensional, commutative)
with $\Lie$ ($(n-1)!$-dimensional, non-commutative) on the closed
colour, while preserving $\Ass$ on the open colour.  This
asymmetry is the operadic origin of the structural difference
between the bulk and the line side: the bulk is governed by
$\Com$ (factorisation on $\Ran(X)$ is symmetric), while the
line-side Koszul dual is governed by $\Lie$ (the Sklyanin
bracket, the Yangian Lie structure).  The open colour, governed
by $\Ass$ on both sides, is the shared associative structure:
the $\Eone$-algebra on the boundary equals the $\Eone$-algebra
on the line side.
\end{remark}


%%%%%%%%%%%%%%%%%%%%%%%%%%%%%%%%%%%%%%%%%%%%%%%%%%%%%%%%%%%%%
\section{PVA descent: $\Ainf$ operations from $\FM_n(\CC)$}
\label{sec:pva-descent}
%%%%%%%%%%%%%%%%%%%%%%%%%%%%%%%%%%%%%%%%%%%%%%%%%%%%%%%%%%%%%

\subsection{The $\Ainf$ operations}

The bar differential $D_\cA$ on
$\barB(\cA) = T^c(s^{-1}\barA)$ is a coderivation of the cofree
coalgebra.  By the universal property of cofree coalgebras, a
coderivation is determined by its composition with the
cogenerator projection
$\mathrm{pr} \colon T^c(s^{-1}\barA) \to s^{-1}\barA$.
The resulting sequence of multilinear maps
\begin{equation}\label{eq:Ainf-ops}
  m_k \colon \cA^{\otimes k}
  \;\longrightarrow\;
  \cA((\lambda_1)) \cdots ((\lambda_{k-1})),
  \qquad k \geq 1, \quad |m_k| = 2-k,
\end{equation}
are the $\Ainf$ operations.  Here $\lambda_j = z_j - z_k$ are
spectral parameters recording the holomorphic separations of
the insertion points on $X$.

The identity $D_\cA^2 = 0$ expands into the
\emph{Stasheff relations}: for each $n \geq 1$,
\begin{equation}\label{eq:stasheff-general}
  \sum_{\substack{i+j = n+1 \\ i,j \geq 1}}
  \;\sum_{s=0}^{n-j}
  (-1)^{\epsilon(s,j)}\;
  m_i\bigl(a_1, \ldots, a_s,\,
  m_j(a_{s+1}, \ldots, a_{s+j}),\,
  a_{s+j+1}, \ldots, a_n\bigr)
  \;=\; 0,
\end{equation}
where $\epsilon(s,j) = \sum_{t=1}^{s}(|a_t|-1)$ is the
Koszul sign from commuting $m_j$ past the desuspended elements.

\subsection{Stasheff relations as Stokes' theorem}

Each Stasheff relation at degree $n$ is Stokes' theorem on the
Fulton--MacPherson space $\FM_n(\CC)$.  The boundary
$\partial\FM_n(\CC)$ is stratified by collision patterns: each
codimension-one stratum $D_S$ (indexed by a subset
$S \subseteq \{1,\ldots,n\}$ with $|S| \geq 2$) records the
collision of the points indexed by $S$.  The contribution of
$D_S$ to the Stasheff relation is $m_{n-|S|+1}$ composed with
$m_{|S|}$, where $m_{|S|}$ acts on the colliding points and
$m_{n-|S|+1}$ acts on the result together with the remaining
points.

\begin{proposition}[Stasheff = Stokes]
\label{prop:stasheff-stokes}
\ClaimStatusProvedHere{}
The $n$-th Stasheff relation \eqref{eq:stasheff-general} is
the identity
\[
  \int_{\FM_n(\CC)} d\omega_n \;=\;
  \int_{\partial\FM_n(\CC)} \omega_n
  \;=\;
  \sum_{|S| \geq 2}
  \int_{D_S} \omega_n,
\]
where $\omega_n$ is the chain-level integrand of $m_n$ contracted
with the weight forms
$\eta_{ij} = d\log(z_i - z_j)$ and OPE coefficients.  The
identity $\partial^2 = 0$ on $\FM_n(\CC)$ (each
codimension-two face is reached by exactly two codimension-one
faces with opposite orientations) ensures $D_\cA^2 = 0$.
\end{proposition}

\begin{proof}
The bar differential is defined as a sum of boundary integrals
on $\FM_n(\CC)$.  The map $m_n$ at degree $n$ is the integral
over the top-dimensional stratum of $\FM_n(\CC)$;
$m_i \circ m_j$ contributions arise from boundary strata.
Stokes' theorem identifies the sum of boundary contributions
with the differential of the bulk integral.  The identity
$d^2 = 0$ on differential forms, applied to the FM
compactification (where boundary strata compose correctly),
gives $D_\cA^2 = 0$.
\end{proof}

We spell out the first four cases.

\subsubsection{$n = 1$: the BRST differential}

$\FM_1(\CC)$ is a single point with empty boundary.  The
Stasheff relation is $m_1 \circ m_1 = 0$: the map $m_1 = Q$
is a differential, $Q^2 = 0$.  This is the BRST condition.

\subsubsection{$n = 2$: the Leibniz rule}

$\FM_2(\CC)$, modulo translations, is a point; its boundary is
the single collision stratum $z_1 = z_2$.  The Stasheff relation
is
\[
  m_1(m_2(a,b))
  \;+\;
  m_2(m_1(a), b)
  \;+\;
  (-1)^{|a|}\, m_2(a, m_1(b))
  \;=\; 0:
\]
the binary product $m_2$ is a chain map with respect to $m_1$.

\subsubsection{$n = 3$: homotopy associativity}

The Stasheff associahedron $K_3$ is a closed interval $[0,1]$,
the quotient of $\FM_3(\CC)$ by the affine group.  Its two
endpoints are the collision strata $D_{12} = \{z_1 \to z_2\}$
and $D_{23} = \{z_2 \to z_3\}$.  The Stasheff relation is:
\begin{equation}\label{eq:stasheff-3}
  m_2(m_2(a,b),c) - m_2(a, m_2(b,c))
  \;=\;
  m_1(m_3(a,b,c)) + m_3(m_1(a),b,c)
  + (-1)^{|a|}\, m_3(a, m_1(b), c)
  + (-1)^{|a|+|b|}\, m_3(a, b, m_1(c)).
\end{equation}
The left-hand side is the associativity defect of $m_2$.  The
right-hand side is the boundary of the homotopy $m_3$: the binary
product is associative up to the homotopy $m_3$, which is the
integral over the interval $K_3$.

When $m_3 = 0$ (Heisenberg, free fermion), the relation reduces to
strict associativity.  When $m_3 \neq 0$ (affine Kac--Moody,
Virasoro), the homotopy is the Jacobiator on $K_3$.

\subsubsection{$n = 4$: the Stasheff pentagon}

The associahedron $K_4$ is a compact polygon with five boundary
faces, corresponding to the collision strata
$D_{12}$, $D_{23}$, $D_{34}$, $D_{123}$, $D_{234}$.  The
Stasheff relation involves $m_2$, $m_3$, and $m_4$, with $m_4$
as the integral over the interior of $K_4$.  For $\cW_3$, the
operation $m_4$ is non-zero, with the coefficient $32/(22+5c)$
entering through the $WW\Lambda$ channel
(Section~\ref{sec:computations}, third computation).


\subsection{PVA descent on cohomology}

\begin{definition}[Poisson vertex algebra]
\label{def:PVA}
A \emph{Poisson vertex algebra} (PVA) over $\Bbbk$ is a
commutative differential algebra $(V, \cdot, \partial)$ equipped
with a $\lambda$-bracket
$\{a_\lambda b\} \in V[\lambda]$
satisfying the following five axioms:
\begin{enumerate}[label=(PVA\arabic*)]
\item \textbf{Sesquilinearity.}
  $\{\partial a_\lambda b\} = -\lambda\{a_\lambda b\}$
  and
  $\{a_\lambda \partial b\} = (\lambda + \partial)\{a_\lambda b\}$.

\item \textbf{Skew-symmetry.}
  $\{a_\lambda b\} = -\{b_{-\lambda-\partial} a\}$.

\item \textbf{Jacobi identity.}
  $\{a_\lambda \{b_\mu c\}\}
  - \{b_\mu \{a_\lambda c\}\}
  = \{\{a_\lambda b\}_{\lambda+\mu} c\}$.

\item \textbf{Left Leibniz rule.}
  $\{a_\lambda bc\}
  = \{a_\lambda b\}c + b\{a_\lambda c\}$.

\item \textbf{Commutativity.}
  $a \cdot b = b \cdot a$.
\end{enumerate}
\end{definition}

\begin{theorem}[PVA descent]
\label{thm:pva-descent}
\ClaimStatusProvedHere{}
Let $\cA$ be a chiral algebra on a smooth curve $X$, and let
$H = H^\bullet(\cA, m_1)$ be the cohomology with respect to
the internal differential $m_1 = Q$.  Suppose
$(H, p, \iota, h)$ is a deformation retract of $(\cA, m_1)$
onto $H$, with the side conditions $h^2 = h\iota = ph = 0$.

Then:
\begin{enumerate}[label=(\alph*)]
\item The homotopy transfer of the $\Ainf$ structure
  $\{m_k\}$ to $H$ produces transferred operations
  $\{m_k^H\}_{k \geq 1}$ satisfying the Stasheff relations.

\item The transferred binary operation $m_2^H$ decomposes as
  \[
    m_2^H(a,b;\lambda)
    \;=\;
    a \cdot b + \{a_\lambda b\},
  \]
  where $a \cdot b = m_2^H(a,b;0)$ is a commutative associative
  product on $H$ and $\{a_\lambda b\}$ is a $\lambda$-bracket.

\item The pair $(H, \cdot, \{-_\lambda -\})$ is a Poisson vertex
  algebra in the sense of Definition~\textup{\ref{def:PVA}}.

\item The higher transferred operations $m_k^H$ for $k \geq 3$
  encode the homotopy corrections to the PVA structure; they
  vanish if and only if $\cA$ is formal (class~G in the shadow
  depth classification).
\end{enumerate}
\end{theorem}

\begin{proof}[Proof outline]
\textbf{Part~(a)} follows from the homotopy transfer theorem
for $\Ainf$ algebras (Kadeishvili \cite{Kad82}, Merkulov
\cite{Mer99}), applied through the homotopy-Koszulity of
$\SCchtop$ (Theorem~\ref{thm:homotopy-koszul-SC}).

\textbf{Part~(b).}  The binary operation $m_2$ has two sources:
the constant term in $\lambda$ (the normally-ordered product,
which is commutative on cohomology) and the polynomial part in
$\lambda$ (the $\lambda$-bracket, which records the OPE
singular part).  On cohomology, the constant term survives as a
commutative product because the $\Sigma_n$-equivariance of the
chiral algebra (for $\Einf$-chiral algebras, all standard vertex
algebras) forces $m_2^H(a,b;0) = m_2^H(b,a;0)$.

\textbf{Part~(c).}  The five PVA axioms are the cohomological
shadows of the Stasheff relations:
\begin{itemize}[nosep]
\item Sesquilinearity (PVA1) is the $n=2$ Stasheff relation
  applied to $\partial a$ and $a$, using the translation
  covariance of the chiral algebra.
\item Skew-symmetry (PVA2) follows from the $\Sigma_2$-equivariance
  of $\FM_2(\CC)$ (the antipodal map $z \mapsto -z$ on the
  boundary circle).
\item The Jacobi identity (PVA3) is the $n=3$ Stasheff relation
  \eqref{eq:stasheff-3} on cohomology, where $m_3$ disappears
  (it is exact on $H$) and the associativity defect becomes the
  Jacobiator.
\item The Leibniz rule (PVA4) is the compatibility of $m_2$
  with the product structure, inherited from the factorisation
  property of the $D$-module.
\item Commutativity (PVA5) is the cohomological consequence of
  $\Einf$-locality, as in part~(b).
\end{itemize}

\textbf{Part~(d).}  If $\cA$ is formal (all $m_k = 0$ for
$k \geq 3$), then $m_2$ is strictly associative and the PVA
structure on $H$ is the complete algebraic datum.  If $\cA$ is
non-formal (class~L, C, or M), the higher operations $m_k^H$
carry the residual $\Ainf$ data that the PVA structure does not
capture.
\end{proof}

\subsection{The five PVA axioms from geometry}
\label{subsec:five-axioms-geometry}

We record the geometric origin of each PVA axiom in detail.

\begin{proposition}[Geometric origin of PVA axioms]
\label{prop:pva-axioms-geometry}
\ClaimStatusProvedHere{}
The five PVA axioms correspond to geometric properties of the
Fulton--MacPherson spaces as follows:
\begin{center}
\renewcommand{\arraystretch}{1.3}
\begin{tabular}{lll}
  \toprule
  \textbf{Axiom} & \textbf{FM property}
  & \textbf{Stasheff} \\
  \midrule
  Sesquilinearity & Translation covariance of $\FM_2(\CC)$
  & $n = 2$ \\
  Skew-symmetry & $\Sigma_2$-action on $\FM_2(\CC)$
  & $n = 2$ \\
  Jacobi & Arnold relation on $\FM_3(\CC)$
  & $n = 3$ \\
  Leibniz & Factorisation on $\FM_2(\CC) \times \FM_1(\CC)$
  & $n = 2$ \\
  Commutativity & $\Einf$-locality ($\Sigma_n$-equivariance)
  & all $n$ \\
  \bottomrule
\end{tabular}
\end{center}
\end{proposition}

\begin{proof}
\textbf{Sesquilinearity.}
The translation $z \mapsto z + c$ acts on $\FM_n(\CC)$ and
preserves all boundary strata.  On the bar complex, translation
covariance is encoded by
$D_\cA \circ \partial = \partial \circ D_\cA$, where $\partial$
is the translation operator on $\cA$.  Applied to a degree-two
bar element $[s^{-1}\partial a \mid s^{-1}b]$, the Stasheff
relation gives
$m_2(\partial a, b; \lambda) = -\lambda\, m_2(a, b; \lambda)$
after extracting the $\lambda$-dependent part (the spectral
parameter $\lambda$ is conjugate to translation, so
$\partial \leftrightarrow -\lambda$ under the Fourier-type
correspondence between position and spectral variables).
The right sesquilinearity follows by skew-symmetry.

\textbf{Skew-symmetry.}
The involution $(z_1, z_2) \mapsto (z_2, z_1)$ on $\FM_2(\CC)$
acts on the boundary circle $E_{12} \cong S^1$ as the antipodal
map.  On the bar complex, this exchanges
$[s^{-1}a \mid s^{-1}b]$ with
$(-1)^{(|a|-1)(|b|-1)}[s^{-1}b \mid s^{-1}a]$.
On the $\lambda$-bracket, the antipodal map sends $\lambda$ to
$-\lambda - \partial$ (the shift by $\partial$ arises from the
chain rule in the Fourier transform), yielding
$\{a_\lambda b\} = -\{b_{-\lambda-\partial} a\}$.

\textbf{Jacobi identity.}
The Arnold relation on $\FM_3(\CC)$,
\begin{equation}\label{eq:arnold-relation}
  \eta_{12} \wedge \eta_{23}
  + \eta_{23} \wedge \eta_{31}
  + \eta_{31} \wedge \eta_{12}
  \;=\; 0,
\end{equation}
where $\eta_{ij} = d\log(z_i - z_j)$, is the fundamental
cohomological relation of the ordered configuration space.
The $n = 3$ Stasheff relation on cohomology (where $m_3$ vanishes)
reduces to
\[
  m_2(m_2(a,b),c) - m_2(a, m_2(b,c)) = 0 \quad (\text{modulo }
  m_3\text{-exact terms}).
\]
Expanding $m_2$ into its product and $\lambda$-bracket components,
the $\lambda$-bracket part of this relation is exactly the
PVA Jacobi identity.  The Arnold relation ensures the three
cyclic terms cancel; this is the same mechanism that gives
$D_\cA^2 = 0$ at degree three.

\textbf{Leibniz rule.}
The factorisation property of the chiral algebra gives a
compatibility between the binary chiral product and the
normally-ordered product.  On $\FM_2(\CC) \times \FM_1(\CC)$
(two colliding points and one spectator), the boundary stratum
where the spectator approaches one of the colliding pair
decomposes the interaction into two terms: $\{a_\lambda b\}c$
and $b\{a_\lambda c\}$.  This is the Leibniz rule.

\textbf{Commutativity.}
For $\Einf$-chiral algebras (all standard vertex algebras), the
chiral product is $\Sigma_n$-equivariant: the factorisation
$D$-module on $\Ran(X)$ is defined on \emph{unordered}
configurations.  On cohomology, this forces the normally-ordered
product to be commutative: $a \cdot b = b \cdot a$.  For
genuinely $\Eone$-chiral algebras (Yangians,
Etingof--Kazhdan quantum vertex algebras), commutativity fails
and the PVA axiom must be replaced by a noncommutative variant.
\end{proof}


%%%%%%%%%%%%%%%%%%%%%%%%%%%%%%%%%%%%%%%%%%%%%%%%%%%%%%%%%%%%%
\section{The resolvent principle}
\label{sec:resolvent}
%%%%%%%%%%%%%%%%%%%%%%%%%%%%%%%%%%%%%%%%%%%%%%%%%%%%%%%%%%%%%

The $\Ainf$ operations $\{m_k\}_{k \geq 2}$ are not independent.
They are determined by a single datum: the binary product $m_2$.

\subsection{The tree-level bootstrap equation}

Let $(H, p, \iota, h)$ be a deformation retract of
$(\cA, m_1)$ onto cohomology $H = H^\bullet(\cA, m_1)$, with
$h^2 = h\iota = ph = 0$.

\begin{definition}[Dressed propagator]
\label{def:dressed-propagator}
The \emph{dressed propagator}
$\Phi \colon T^+(H) \to \cA$ is defined by the fixed-point
equation
\begin{equation}\label{eq:dressed-propagator}
  \Phi
  \;=\;
  \iota + h \circ m_2 \circ \Phi^{\otimes 2}.
\end{equation}
Here $T^+(H) = \bigoplus_{k \geq 1} H^{\otimes k}$ is the
reduced tensor module, $\iota \colon H \hookrightarrow \cA$ is
the inclusion, and $h \colon \cA \to \cA$ is the contracting
homotopy.
\end{definition}

Equation~\eqref{eq:dressed-propagator} is the tree-level
bootstrap: the full propagator from boundary data $H$ into the
bulk algebra $\cA$ is determined self-consistently from the free
embedding $\iota$ and the interaction vertex $m_2$.  Iteration
produces the tree formula.

\subsection{The tree formula}

\begin{theorem}[Resolvent principle]
\label{thm:resolvent}
\ClaimStatusProvedHere{}
The transferred $\Ainf$ operations on $H$ are given by
\begin{equation}\label{eq:tree-formula}
  m_k^H
  \;=\;
  \sum_{T \in \mathrm{PRT}_k}
  p \circ \mathfrak{m}(T) \circ \iota^{\otimes k},
\end{equation}
where:
\begin{enumerate}[label=(\roman*)]
\item $\mathrm{PRT}_k$ is the set of planar rooted binary trees
  with $k$ leaves;
\item $|\mathrm{PRT}_k| = C_{k-1}$, the $(k-1)$-st Catalan number
  ($C_0 = 1$, $C_1 = 1$, $C_2 = 2$, $C_3 = 5$, \ldots);
\item $\mathfrak{m}(T)$ assigns the binary product $m_2$ to each
  internal vertex and the contracting homotopy $h$ to each
  internal edge;
\item $p \colon \cA \to H$ is the projection and
  $\iota \colon H \to \cA$ is the inclusion.
\end{enumerate}
\end{theorem}

\begin{proof}
The iteration of \eqref{eq:dressed-propagator} at order $k$ collects
all ways of building a $k$-to-$1$ map from $H^{\otimes k}$ to
$\cA$ using the vertex $m_2$ and the propagator $h$.  Each such
composition pattern is a planar rooted binary tree with $k$
leaves (one leaf per input from $H$), internal vertices labelled
by $m_2$, and internal edges labelled by $h$.  The projection $p$
extracts the cohomology class.

The number of planar rooted binary trees with $k$ leaves is the
Catalan number $C_{k-1}$: this follows from the recursive
decomposition of a tree into its left and right subtrees,
which gives the recurrence
$C_n = \sum_{i=0}^{n-1} C_i C_{n-1-i}$ with $C_0 = 1$.

The Stasheff relations for $\{m_k^H\}$ are verified by
checking that the boundary contributions of each tree (where
two adjacent vertices merge) produce the correct
composition terms.  This is the combinatorial content of the
associahedron: the faces of $K_n$ correspond to the degenerate
trees.
\end{proof}

One equation, one binary input, $C_{k-1}$ trees.

\begin{remark}[Physical interpretation]
\label{rem:physical-tree-formula}
The tree formula is the perturbative expansion of the Lagrangian
self-intersection $\cL \times_{\cM}^h \cL$ in a shifted symplectic
stack.  Each tree is a Feynman diagram with vertices $m_2$ and
propagators $h$.  The sum over $\mathrm{PRT}_k$ is the tree-level
scattering amplitude at multiplicity $k$.  Loop corrections
(genus $g \geq 1$) arise from the curvature
$d_{\mathrm{fib}}^2 = \kappa(\cA) \cdot \omega_g$, which obstructs
the tree-level bootstrap at higher genus.
\end{remark}


%%%%%%%%%%%%%%%%%%%%%%%%%%%%%%%%%%%%%%%%%%%%%%%%%%%%%%%%%%%%%
\section{Three computations}
\label{sec:computations}
%%%%%%%%%%%%%%%%%%%%%%%%%%%%%%%%%%%%%%%%%%%%%%%%%%%%%%%%%%%%%

\subsection{First computation: Heisenberg (formal, class~G)}
\label{subsec:comp-heisenberg}

Let $\cA = \cH_k$ be the Heisenberg algebra at level $k$, with
generator $J$ and OPE
\[
  J(z)\, J(w)
  \;\sim\;
  \frac{k}{(z-w)^2}.
\]

\begin{proposition}[Heisenberg $\Ainf$ structure]
\label{prop:heisenberg-ainf}
\ClaimStatusProvedHere{}
The $\Ainf$ structure of $\cH_k$ is:
\begin{enumerate}[label=(\roman*)]
\item $m_1 = 0$ (no internal differential).
\item $m_2(J, J; \lambda) = k\lambda$ (the double-pole residue,
  absorbed by $d\log$: the OPE has a double pole, and the
  $d\log(z_1 - z_2)$ kernel reduces the pole order by one,
  producing a simple-pole collision residue $k$ times the spectral
  parameter).
\item $m_k = 0$ for all $k \geq 3$ (no composites, no higher
  poles).
\end{enumerate}
The algebra is formal: the $\Ainf$ structure is concentrated in
$m_2$.
\end{proposition}

\begin{proof}
The augmentation-ideal projected binary product
$\bar\mu_2(J, J) = 0$: the OPE $J(z)J(w) \sim k/(z-w)^2$ has
singular part a scalar (the identity), which projects to
zero in $\barA = \ker(\varepsilon)$.  Therefore the bar
differential on $\barB(\cH_k)$ vanishes identically: all bar
elements are $D_\cA$-closed.

The operation $m_2(J,J;\lambda) = k\lambda$ is obtained from
the full OPE by the $d\log$ extraction: the coefficient of
$1/(z_1 - z_2)^2$ in the OPE is $k$, and the spectral
parameter $\lambda = z_1 - z_2$ after Fourier transform gives
$k\lambda$.

Since $\bar\mu_2 = 0$ on $\barA$, the degree-$n$ bar
differential involves only the identity-valued OPE residue,
which lives outside $\barA$.  All $m_k$ for $k \geq 3$ vanish.
\end{proof}

\begin{corollary}[Heisenberg PVA]
\label{cor:heisenberg-pva}
On $H^\bullet(\cH_k) = \Bbbk[J]$, the PVA structure is:
\begin{enumerate}[label=(\roman*)]
\item Product: $J^n \cdot J^m = J^{n+m}$ (polynomial ring).
\item $\lambda$-bracket: $\{J_\lambda J\} = k\lambda$.
\item All higher operations $m_k^H = 0$ for $k \geq 3$.
\end{enumerate}
This is a PVA of rank one.  The classical $r$-matrix is
$r(z) = k/z$.

Verification: at $k = 0$, the $r$-matrix vanishes ($r(z) = 0$)
and the algebra reduces to a commutative ring with trivial
$\lambda$-bracket: the abelian limit.

The modular characteristic is $\kappa(\cH_k) = k$: the same
scalar governs the OPE, the $\lambda$-bracket, the $r$-matrix,
and the genus-one curvature $d_{\mathrm{fib}}^2 = k \cdot \omega_1$.
\end{corollary}

\begin{remark}[Shadow depth zero]
\label{rem:heisenberg-shadow}
The Heisenberg algebra belongs to class~G in the shadow depth
classification: $d_{\mathrm{alg}} = 0$.  The shadow tower
terminates at $\kappa = k$; the discriminant
$\Delta = 8\kappa S_4 = 0$ (since $S_4 = 0$ for the Heisenberg)
confirms the finite tower.  The SC-formality condition is
satisfied: $\cH_k$ is SC-formal because $m_k = 0$ for $k \geq 3$.
\end{remark}


\subsection{Second computation: Kac--Moody at level $k$
  (Lie-transverse, class~L)}
\label{subsec:comp-km}

Let $\cA = V_k(\fg)$ be the vacuum module of the affine
Kac--Moody algebra at level
$k \in \CC \setminus \{-h^\vee\}$, with generators $\{J^a\}$
($a = 1, \ldots, d = \dim\fg$) and OPE
\begin{equation}\label{eq:km-ope}
  J^a(z)\, J^b(w)
  \;\sim\;
  \frac{k\, \kappa^{ab}}{(z-w)^2}
  \;+\;
  \frac{f^{ab}{}_{c}\, J^c(w)}{z-w},
\end{equation}
where $\kappa^{ab}$ is the normalised Killing form and
$f^{ab}{}_c$ are the structure constants.

\begin{proposition}[Kac--Moody $\Ainf$ structure]
\label{prop:km-ainf}
\ClaimStatusProvedHere{}
The $\Ainf$ structure of $V_k(\fg)$ is:
\begin{enumerate}[label=(\roman*)]
\item $m_1 = 0$.
\item $m_2(J^a, J^b; \lambda) = f^{ab}{}_c\, J^c + k\,\kappa^{ab}\,\lambda$.
  The simple-pole coefficient $f^{ab}{}_c J^c$ is the Lie bracket;
  the double-pole coefficient $k\,\kappa^{ab}$ contributes the
  $\lambda$-dependent term.
\item $m_3 \neq 0$: the Jacobiator homotopy on $\FM_3(\CC)$.
  The three codimension-one boundary strata of $\FM_3(\CC)$
  contribute $m_2(m_2(-,-),-)$ terms; their failure to cancel
  (the Jacobi identity defect from the simple-pole bracket) is
  corrected by $m_3$.
\item $m_k = 0$ for all $k \geq 4$, by dimension counting on
  $\FM_k(\CC)$: the OPE has at most a double pole, so the
  chain-level integrands vanish for $k \geq 4$.
\end{enumerate}
\end{proposition}

\begin{proof}
The binary operation extracts from the OPE
\eqref{eq:km-ope}: the residue at $z_1 = z_2$ with the
$d\log(z_1 - z_2)$ kernel gives the simple-pole coefficient
$f^{ab}{}_c J^c$ (the Lie bracket), and the double-pole
coefficient $k\kappa^{ab}$ contributes the $\lambda$-linear
term.

The existence of $m_3 \neq 0$ follows from the failure of
strict associativity: $m_2(m_2(J^a, J^b), J^c) -
m_2(J^a, m_2(J^b, J^c)) = f^{ab}{}_e f^{ec}{}_d J^d -
f^{bc}{}_e f^{ae}{}_d J^d \neq 0$ (the Jacobi defect).  The
homotopy $m_3$ is the integral over $K_3 = [0,1]$ that
interpolates between the two boundary contributions.

The vanishing $m_k = 0$ for $k \geq 4$ follows from the pole
structure: the OPE has at most double poles, and the weight
forms $\eta_{ij}$ absorb one pole order each.  For $k \geq 4$
insertion points, the integrand on $\FM_k(\CC)$ involves at most
$(k-1)$ weight forms but the pole excess is at most $2$, which
is consumed at $k = 3$.  For $k \geq 4$, the integrand has no
residual singularity and the FM integral vanishes.
\end{proof}

\begin{corollary}[Kac--Moody PVA]
\label{cor:km-pva}
On $H^\bullet(V_k(\fg)) = \Bbbk[J^a : a = 1, \ldots, d]$, the
PVA structure is:
\begin{enumerate}[label=(\roman*)]
\item Product: polynomial ring in $d$ generators.
\item $\lambda$-bracket:
  $\{J^a_\lambda J^b\} = f^{ab}{}_c\, J^c + k\,\kappa^{ab}\,\lambda$.
\item Classical $r$-matrix: $r(z) = k\,\Omega/z$ in the
  trace-form convention, where $\Omega = \sum_a J^a \otimes J_a$ is
  the split Casimir.

  Verification: at $k = 0$, the $r$-matrix vanishes
  ($r(z) = 0$); the algebra reduces to the commutative ring with
  Lie bracket $\{J^a_\lambda J^b\} = f^{ab}{}_c J^c$, which is
  still a well-defined PVA (the abelian limit of the level).

  In the KZ convention: $r(z) = \Omega/((k+h^\vee)z)$.  Bridge:
  $k\,\Omega_{\mathrm{tr}} = \Omega/(k+h^\vee)$ at generic $k$.
\end{enumerate}

The modular characteristic is
\begin{equation}\label{eq:kappa-km}
  \kappa(V_k(\fg))
  \;=\;
  \frac{\dim(\fg)\,(k + h^\vee)}{2h^\vee}.
\end{equation}
At $k = 0$: $\kappa = \dim(\fg)/2$ (not zero: the pure vacuum
anomaly).  At $k = -h^\vee$: $\kappa = 0$ (critical level, the
bar complex is flat).
\end{corollary}

\begin{remark}[Shadow depth one, class~L]
\label{rem:km-shadow}
The affine Kac--Moody algebra belongs to class~L:
$d_{\mathrm{alg}} = 1$.  The shadow tower extends to the cubic
invariant $\mathfrak{C}(\hat\fg_k)$ (from the Lie structure
constants) but terminates at shadow depth $r_{\max} = 3$.  The
quartic shadow invariant $S_4 = 0$, so the discriminant
$\Delta = 8\kappa S_4 = 0$: the shadow tower is finite
(class~L terminates at depth $3$).  The algebra is
\emph{not} SC-formal:
$m_3 \neq 0$.
\end{remark}


\subsection{Third computation: $\cW_3$ (infinite tower, class~M)}
\label{subsec:comp-w3}

The $\cW_3$ algebra has two generators: the stress tensor $T$
of conformal weight $2$ and a spin-$3$ current $W$ of conformal
weight $3$.

\begin{proposition}[$\cW_3$ $\Ainf$ structure]
\label{prop:w3-ainf}
\ClaimStatusProvedHere{}
The $\Ainf$ structure of $\cW_3$ satisfies:
\begin{enumerate}[label=(\roman*)]
\item $m_1 = 0$.
\item $m_2(T, T; \lambda) = (c/12)\lambda^3 + 2T\lambda + \partial T$
  (the full $TT$ OPE in divided-power parametrisation; the
  coefficient $c/12$ matches $T_{(3)}T = c/2$ via the divided
  power $\lambda^{(3)} = \lambda^3/3!$).
\item $m_2(T, W; \lambda) = 3W\lambda + \partial W$.
\item The $WW$ operation involves the composite field
  $\Lambda = {}{:}(TT){:} - \frac{3}{10}\partial^2 T$ at the
  $\lambda$-linear term:
  \[
    m_2(W, W; \lambda)
    \;=\;
    \frac{c}{360}\lambda^5
    + \frac{1}{3}T\lambda^3
    + \frac{1}{2}(\partial T)\lambda^2
    + \Bigl(\frac{3}{10}\partial^2 T
      + \frac{32}{22+5c}\,\Lambda\Bigr)\lambda
    + \cdots
  \]
  The coefficient $32/(22+5c)$ diverges at $c = -22/5$
  (the Lee--Yang value).
\item $m_3 \neq 0$: the Jacobiator homotopy, involving $\Lambda$.
\item $m_4 \neq 0$: the quartic operation from the $\FM_4(\CC)$
  boundary decomposition, with $32/(22+5c)$ entering through the
  $WW\Lambda$ channel.
\item The $\Ainf$ structure does not terminate: $m_k \neq 0$
  for all $k \geq 3$.
\end{enumerate}
\end{proposition}

\begin{proof}[Proof sketch]
The binary operations are extracted from the Zamolodchikov
$\cW_3$ OPE \cite{Zam85} using the $d\log$ kernel and divided-power
parametrisation.  The non-vanishing of $m_3$ follows from the
non-Lie structure of the $\cW_3$ bracket (the bracket involves
the composite field $\Lambda$, which is not a generator).  The
non-vanishing of $m_4$ follows from the Stasheff pentagon: the
$m_3$-$m_3$ composition on $\partial K_4$ does not cancel (the
$\Lambda$ channel provides a non-trivial residue on the
interior of $K_4$).  The infinite tower $m_k \neq 0$ for all $k$
follows by induction: each $m_k$ generates a non-vanishing
correction at degree $k+1$ through the Stasheff cascade, driven
by the non-primitive composite $\Lambda$.
\end{proof}

\begin{corollary}[$\cW_3$ PVA]
\label{cor:w3-pva}
On cohomology, the PVA structure is the Zamolodchikov PVA:
\begin{align*}
  \{T_\lambda T\}
  &\;=\;
  \frac{c}{12}\lambda^3 + 2T\lambda + \partial T, \\
  \{T_\lambda W\}
  &\;=\;
  3W\lambda + \partial W, \\
  \{W_\lambda W\}
  &\;=\;
  \frac{c}{360}\lambda^5
  + \frac{1}{3}T\lambda^3
  + \frac{1}{2}(\partial T)\lambda^2
  + \Bigl(\frac{3}{10}\partial^2 T
    + \frac{32}{22+5c}\,\Lambda\Bigr)\lambda
  + \cdots
\end{align*}
The modular characteristic is $\kappa(\cW_3) = 5c/6$: from the
formula $\kappa(\cW_N) = c\,(H_N - 1)$ with
$H_N = \sum_{j=1}^{N} 1/j$, at $N = 3$ we have
$H_3 = 1 + 1/2 + 1/3 = 11/6$, so $\kappa = c \cdot 5/6$.

Verification at $N = 2$: $H_2 = 3/2$, $H_2 - 1 = 1/2$, so
$\kappa(\cW_2) = c/2$, which matches the Virasoro formula
$\kappa(\mathrm{Vir}_c) = c/2$ (since $\cW_2 = \mathrm{Vir}$).
\end{corollary}

\begin{remark}[Class~M: infinite shadow depth]
\label{rem:w3-shadow}
The $\cW_3$ algebra belongs to class~M:
$d_{\mathrm{alg}} = \infty$.  The shadow tower
$S_2, S_3, S_4, \ldots$ is infinite, driven by the
non-primitive composite $\Lambda$.  The generating depth
$d_{\mathrm{gen}} = 3$ (the operation $m_3$ generates all higher
$m_k$ via the Stasheff cascade), but the algebraic depth is
infinite (all $m_k \neq 0$).  The discriminant
$\Delta = 8\kappa S_4 \neq 0$ for $c \neq 0$ and
$c \neq -22/5$, confirming the infinite tower.

This infinite depth is the algebraic shadow of the fact that
three-dimensional gravity with $\cW_3$ symmetry requires an
infinite number of graviton vertices: the $\Ainf$ structure is
the perturbative expansion of the bulk theory.
\end{remark}

\begin{center}
\renewcommand{\arraystretch}{1.3}
\begin{tabular}{lccccl}
  \toprule
  \textbf{Algebra} & \textbf{Class} & $d_{\mathrm{alg}}$ &
  $\kappa$ & $m_{k \geq 3}$ & \textbf{PVA type} \\
  \midrule
  $\cH_k$ & G & $0$ & $k$ & $= 0$ & free rank-1 \\
  $V_k(\fg)$ & L & $1$ & $\frac{d(k+h^\vee)}{2h^\vee}$ &
  $m_3 \neq 0$, $m_{\geq 4} = 0$ & affine, Lie-transverse \\
  $\cW_3$ & M & $\infty$ & $5c/6$ & all $\neq 0$ &
  Zamolodchikov, infinite \\
  \bottomrule
\end{tabular}
\end{center}


%%%%%%%%%%%%%%%%%%%%%%%%%%%%%%%%%%%%%%%%%%%%%%%%%%%%%%%%%%%%%
%% Wave-14 reconception: heptagon embedding, topologization
%%    scope, PVA descent via Khan-Zeng, AP-SC discipline
%%%%%%%%%%%%%%%%%%%%%%%%%%%%%%%%%%%%%%%%%%%%%%%%%%%%%%%%%%%%%

\section{Wave-14 reconception: heptagon, topologization scope,
 PVA descent}
\label{sec:wave14-schtop-heptagon}

This addendum reconceives the operad $\SCchtop$ against three
2026-04-17 platonic-ideal anchors: the Vol~II $\SCchtop$ heptagon
(seven edges routing bar, cobar, Drinfeld centre, derived centre,
Swiss-cheese, PTVV, celestial), the four-class topologisation scope
tree, and the Khan--Zeng PVA descent for freely-generated Poisson
vertex algebras.

\begin{theorem}[The $\SCchtop$ heptagon]
\label{thm:wave14-heptagon}
The two-coloured operad $\SCchtop$ sits at the hub of a heptagon of
seven equivalent presentations, each edge a named theorem:
\begin{enumerate}[(1)]
\item \emph{Operadic hub:} $\SCchtop$ itself (this paper).
\item \emph{Koszul dual:} $(\SCchtop)^! = (\Lie^c, \Ass^c,
  \textup{shuffle-mixed})$ (this paper, Theorem~\ref{thm:SC-koszul-dual};
  \emph{not} Koszul self-dual).
\item \emph{Chiral factorization:} the Beilinson--Drinfeld
  factorization algebra on $\Ran(X)$ with open $\bR$-decoration.
\item \emph{BV--BRST:} the Costello--Gwilliam bulk--boundary
  BV complex.
\item \emph{Convolution:} the $L_\infty$-algebra
  $\Hom_\alpha(C,A)$ with open-closed decoration.
\item \emph{Drinfeld centre:} $Z(\Rep_{\mathrm{fact}}(\cA)) \simeq
  \Rep_{\mathrm{fact}}(Z^{\mathrm{der}}_{\mathrm{ch}}(\cA))^{E_2}$
  via categorified bar--cobar with half-braiding.
\item \emph{PTVV derived-AG:} the $1$-shifted symplectic mapping
  stack $\mathrm{Map}(X \times \bR_{\ge 0}, B\,\SC\text{-}\mathrm{Alg})$.
\end{enumerate}
Each of the seven is a distinct theorem; the Vol~II heptagon
chapter (\texttt{sc\_chtop\_heptagon.tex}) proves the seven
edges as named theorems.
\end{theorem}

\begin{theorem}[Topologization scope tree, AP-TOPOLOGIZATION]
\label{thm:wave14-topologization-scope}
The passage $\SCchtop \to E_3^{\mathrm{top}}$ via Sugawara
topologisation branches by shadow class:
\begin{itemize}
\item Class~G (Heisenberg): $E_3^{\mathrm{top}}$ immediate;
  commutative input, Dunn on $F^{H_k}$.
\item Class~L ($V_k(\fg)$, $k \ne -h^\vee$):
  $E_3^{\mathrm{top}}$ \emph{proved on the original complex}, via
  Costello--Gwilliam factorization on $\bR \times \bC$ and explicit
  BRST primitives $\eta_1^{(i)}, \eta_1^{(ii)}$ giving
  $[Q, \widetilde G_1] = T_{\mathrm{Sug}}$ at strict chain level.
\item Class~C ($\beta\gamma$, $bc$): $E_3^{\mathrm{top}}$ via
  Friedan--Martinec--Shenker bosonisation reducing to class~G.
\item Class~M ($\Vir$, $\cW_N$): $E_3^{\mathrm{top}}$
  \emph{weight-completed}, via half-BRST plus MC5 pattern;
  original-complex chain-level is the sole remaining open
  direction.
\item Critical level $k = -h^\vee$: $E_2^{\mathrm{top}}$
  (a dimension drop, not a failure; Sugawara is undefined at
  critical level).
\end{itemize}
Without a conformal vector (or at critical level), the theory is
stuck at $\SCchtop$; this is the generic situation (Heisenberg,
lattice, $\Eone$-chiral Yangian input, free fields without stress
tensor), for which $\SCchtop$ is the final answer and $E_3$ the
special case.
\end{theorem}

\begin{theorem}[Khan--Zeng PVA descent on freely generated PVAs]
\label{thm:wave14-khan-zeng}
Let $A$ be a freely-generated Poisson vertex algebra with
conformal vector. Then the $\Ainf$ operations $\{m_k\}$ extracted
by cogenerator projection from the bar differential of the
associated graded $\mathrm{gr}_{Li}(A)$ descend to the five PVA
axioms (sesquilinearity, skew-symmetry, Jacobi, Leibniz,
commutativity) on cohomology. The descent covers the
3d Poisson sigma-model class of Khan--Zeng, comprising all
freely-generated PVAs with conformal vector; in particular, it
recovers the Vol~II orbifold route via $\bZ/n$-invariants
preserving $\En$-structure.
\end{theorem}

\begin{remark}[AP-SC-BAR discipline reaffirmed]
\label{rem:wave14-ap-sc-bar}
As stated in Warning~\ref{warn:bar-not-SC}, the bar complex
$\barB(\cA) = T^c(s^{-1}\barA)$ is an $\Eone$ chiral
coassociative coalgebra, \emph{not} an $\SCchtop$-coalgebra. The
$\SCchtop$ structure lives on the derived centre pair
$(C^\bullet_{\mathrm{ch}}(\cA,\cA), \cA)$; the bar complex serves
as a resolution used to \emph{compute}
$C^\bullet_{\mathrm{ch}}(\cA,\cA)$, not as an object carrying
$\SCchtop$-structure. The six-step chain of errors (FM25) is
explicitly retired: (1)~the bar differential is not a closed
colour; (2)~deconcatenation is not an open colour; (3)~two
structures on a single coalgebra do not produce two colours;
(4)~$\SCchtop$ emerges in the output via
$\Hom(B(\cA), \cA)$, not in the input; (5)~$\Eone$-chiral
$\ne$ $\Eone$-topological (chiral lives on a curve, topological
on $\Conf_k(\bR)$); (6)~the Steinberg analogy is retired.
\end{remark}

\begin{remark}[AP-SC-NOT-SELFDUAL reaffirmed]
\label{rem:wave14-ap-sc-notselfdual}
$(\SCchtop)^! \ne \SCchtop$: closed-colour dimensions differ
at $n \ge 3$ ($\dim\Com(n) = 1$ vs $\dim\Lie(n) = (n-1)!$).
The duality \emph{functor} is an involution; the operad
is \emph{not} self-dual. Livernet's proof gives Koszulity, not
self-duality.
\end{remark}

\begin{remark}[Epistemic tags]
\label{rem:wave14-schtop-epistemic}
Theorem~\ref{thm:wave14-heptagon}: the seven heptagon edges are
\emph{proved} in the Vol~II heptagon chapter; this paper's
operadic hub (edges 1--3) is proved directly.
Theorem~\ref{thm:wave14-topologization-scope}: class~G, L, C
are unconditional on the original complex; class~M is
weight-completed (AP-TOPOLOGIZATION). Theorem~\ref{thm:wave14-khan-zeng}:
unconditional for freely-generated PVAs; the general class~M case
reduces to the weight-completed topologization.
\end{remark}


%%%%%%%%%%%%%%%%%%%%%%%%%%%%%%%%%%%%%%%%%%%%%%%%%%%%%%%%%%%%%
%% REFERENCES
%%%%%%%%%%%%%%%%%%%%%%%%%%%%%%%%%%%%%%%%%%%%%%%%%%%%%%%%%%%%%

\begin{thebibliography}{99}

\bibitem{AF15}
D.~Ayala and J.~Francis,
\emph{Factorization homology of topological manifolds},
J.~Topol.~\textbf{8} (2015), no.~4, 1045--1084.

\bibitem{BD04}
A.~Beilinson and V.~Drinfeld,
\emph{Chiral Algebras},
Amer.\ Math.\ Soc.\ Colloq.\ Publ., vol.~51, AMS, 2004.

\bibitem{Cav14}
G.~Caviglia,
\emph{The Quillen model structure for coloured operads},
arXiv:1401.6983, 2014.

\bibitem{DK10}
V.~Dotsenko and A.~Khoroshkin,
\emph{Gr\"obner bases for operads},
Duke Math.~J.~\textbf{153} (2010), no.~2, 363--396.

\bibitem{DSK06}
A.~De~Sole and V.~Kac,
\emph{Finite vs.\ affine $W$-algebras},
Jpn.~J.~Math.~\textbf{1} (2006), no.~1, 137--261.

\bibitem{FG12}
J.~Francis and D.~Gaitsgory,
\emph{Chiral Koszul duality},
Selecta Math.~(N.S.)~\textbf{18} (2012), no.~1, 27--87.

\bibitem{GK94}
V.~Ginzburg and M.~Kapranov,
\emph{Koszul duality for operads},
Duke Math.~J.~\textbf{76} (1994), no.~1, 203--272.

\bibitem{Hin97}
V.~Hinich,
\emph{Homological algebra of homotopy algebras},
Comm.~Algebra~\textbf{25} (1997), no.~10, 3291--3323.

\bibitem{Kad82}
T.~V.~Kadeishvili,
\emph{On the theory of homology of fiber spaces},
Uspekhi Mat.~Nauk~\textbf{35} (1980), no.~3, 183--188;
Int.\ Math.\ Res.\ Not.\ (2003), no.~3, 141.

\bibitem{Kon99}
M.~Kontsevich,
\emph{Operads and motives in deformation quantization},
Lett.~Math.~Phys.~\textbf{48} (1999), no.~1, 35--72.

\bibitem{Kon03}
M.~Kontsevich,
\emph{Deformation quantization of Poisson manifolds},
Lett.~Math.~Phys.~\textbf{66} (2003), no.~3, 157--216.

\bibitem{KS06}
H.~Kajiura and J.~Stasheff,
\emph{Open-closed homotopy algebra in mathematical physics},
J.~Math.~Phys.~\textbf{47} (2006), 023506.

\bibitem{Liv06}
M.~Livernet,
\emph{A rigidity theorem for pre-Lie algebras},
J.~Pure Appl.~Algebra~\textbf{207} (2006), no.~1, 1--18.

\bibitem{LV12}
J.-L.~Loday and B.~Vallette,
\emph{Algebraic Operads},
Grundlehren Math.~Wiss., vol.~346, Springer, 2012.

\bibitem{Mer99}
S.~Merkulov,
\emph{Strong homotopy algebras of a K\"ahler manifold},
Internat.\ Math.\ Res.\ Notices~(1999), no.~3, 153--164.

\bibitem{Vor99}
A.~Voronov,
\emph{The Swiss-cheese operad},
Contemp.~Math.~\textbf{239} (1999), 365--373.

\bibitem{Zam85}
A.~B.~Zamolodchikov,
\emph{Infinite additional symmetries in two-dimensional
conformal quantum field theory},
Theor.~Math.~Phys.~\textbf{65} (1985), no.~3, 1205--1213.

\end{thebibliography}

\end{document}
